\documentclass[11pt,a4paper]{article}

\usepackage[utf8]{inputenc}
\usepackage[T1]{fontenc}
\usepackage{newtxtext}
\usepackage[
    top=1in,
    bottom=1in,
    left=1.2in,
    right=1in
]{geometry}
\usepackage[table]{xcolor}
\usepackage{tabularx}
\usepackage{array}
\usepackage{booktabs}
\usepackage{enumitem}
\usepackage{graphicx}
\usepackage{titlesec}
\usepackage{fancyhdr}
\usepackage{parskip}
\usepackage[hidelinks]{hyperref}
 

\definecolor{srsnavy}{HTML}{1B4332}
\definecolor{srsblue}{HTML}{2D6A4F}
\definecolor{srsgrey}{HTML}{495057}
 
\titleformat{\section}{\Large\bfseries\color{srsnavy}}{\thesection}{0.6em}{}
\titleformat{\subsection}{\large\bfseries\color{srsblue}}{\thesubsection}{0.6em}{}
\titleformat{\subsubsection}{\normalsize\bfseries\color{srsgrey}}{\thesubsubsection}{0.6em}{}
 

\pagestyle{fancy}
\fancyhf{}
\renewcommand{\headrulewidth}{0.4pt}
\renewcommand{\footrulewidth}{0.4pt}
\fancyhead[L]{\small\textcolor{srsgrey}{\labtitle}}
\fancyhead[R]{\small\textcolor{srsgrey}{\coursename}}
\fancyfoot[L]{\small\textcolor{srsgrey}{\institution}}
\fancyfoot[R]{\small\textcolor{srsgrey}{Page \thepage}}
 

\newcommand{\reqheader}[3]{%
  \rowcolor{srsnavy}
  \textcolor{white}{\textbf{#1}} &
  \textcolor{white}{\textbf{#2}} &
  \textcolor{white}{\textbf{#3}} \\ \hline}
 

\newcommand{\institution}{IOE, Purwanchal Campus}      % used in the footer
\newcommand{\department}{Department of Computer \& Electronics Engineering}
\newcommand{\logofile}{asdf.png}
\newcommand{\coursename}{Software Engineering (Practical)}
\newcommand{\coursecode}{ ENCT 352 }
\newcommand{\labtitle}{Lab 1 -- Software Requirements Specification}
\newcommand{\projecttitle}{Library Management System}   % <-- your project
\newcommand{\submissiondate}{\today}
\newcommand{\teamname}{Group XYZ}
%==============================================================

\begin{document}
 
%--------------------------------------------------------------
%  TITLE PAGE
%--------------------------------------------------------------
\begin{titlepage}
\centering
\vspace*{0.6cm}
 
% ---- LOGO (shown if the image file exists; otherwise a placeholder) ----
\IfFileExists{\logofile}{%
  \includegraphics[width=15cm]{\logofile}%
}{%
  \fbox{\parbox[c][3cm][c]{3cm}{\centering\footnotesize Place\\ \texttt{\logofile}\\ beside this file}}%
}\par

\vspace{0.6cm}
{\Huge\bfseries\color{srsnavy} Software Requirements Specification\par}
\vspace{0.3cm}
{\large\itshape Project: \projecttitle\par}
\vspace{0.25cm}
{\department\par}
\vspace{1.5cm}
 
% ---- submission details ----
\renewcommand{\arraystretch}{1.4}
\begin{tabular}{>{\bfseries}l l}
Course        & \coursename \\
Course code   & \coursecode \\
Team          & \teamname \\
Submitted on  & \submissiondate \\
\end{tabular}
 
\vspace{1.2cm}
% ---- team members: add/remove rows ----
{\bfseries\color{srsblue} Team Members}\par
\vspace{0.3cm}
\renewcommand{\arraystretch}{1.3}
\begin{tabular}{|c|l|l|}
\hline
\rowcolor{srsnavy}
\textcolor{white}{\textbf{S.N.}} & \textcolor{white}{\textbf{Name}} & \textcolor{white}{\textbf{Roll / ID}} \\ \hline
1 & ABC  ABC & PUR074BCT045 \\ \hline
2 & ABC  ABC& PUR074BCT045 \\ \hline
3 & ABC  ABC& PUR074BCT045 \\ \hline
4 & ABC  ABC & PUR074BCT045 \\ \hline
\end{tabular}
 
\vfill
{\small\color{srsgrey} Submitted to: Asst prof. Binay Lal Shrestha }
\end{titlepage}

\clearpage
\tableofcontents
\clearpage

%==============================================================
%  1. INTRODUCTION
%==============================================================
\section{Introduction}

\subsection{Purpose}
% State what this SRS is for and who should read it.
This document specifies the requirements for the \projecttitle{}. It is intended
for the development team, the instructor, and any stakeholder who needs to
understand what the system must do before it is built.

\subsection{Scope}
% What the system will do, the benefit it provides, and what is OUT of scope.
The \projecttitle{} will allow librarians to manage a book catalogue and members,
and allow members to search and reserve books. It replaces the existing manual
register and enforces borrowing rules automatically. Out of scope: payment
processing and inter-library transfers.

\subsection{Definitions, Acronyms, and Abbreviations}
\begin{description}[leftmargin=2.5cm,style=nextline]
  \item[SRS] Software Requirements Specification
  \item[LMS] Library Management System
  \item[ISBN] International Standard Book Number
  \item[FR / NFR] Functional / Non-functional Requirement
\end{description}

\subsection{Overview of the Document}
Summarises the SRS.

\section{Objectives}
\textbf{Aim:} To elicit the requirements of the chosen system and document them as a
Software Requirements Specification using the IEEE~830 template.

\textbf{Tools used:} \LaTeX{} for the document.


\section{Project Domain Analysis}

\subsection{Domain Description}
A library maintains a collection of books, registers members, and tracks loans. Typical operations: add new books, register members, issue books, return books, and handle reservations. The domain involves two primary actors: Librarian (staff) and Member (patron).

\subsection{Existing System Overview}
Currently, the library uses a paper‑based register. The librarian manually records book issues, due dates, and returns. Searching for a book requires browsing shelves or card catalogues. Reservation is not possible; members must visit physically to check availability.

\subsection{Problems in Existing System}

The existing library management system is largely manual and relies on paper-based records for maintaining book inventories, member information, and borrowing transactions. This approach presents several challenges:

\begin{itemize}[leftmargin=*]
\item \textbf{Time-consuming operations:} Recording book issues, returns, and member details manually requires significant time and effort.

\item \textbf{Difficulty in searching records:} Locating information about a particular book or member from physical registers is slow and inefficient.

\item \textbf{Human errors:} Manual data entry can lead to mistakes in book records, issue dates, return dates, and member information.

\item \textbf{Poor tracking of book availability:} Determining whether a book is available, issued, or reserved is difficult without real-time updates.

\item \textbf{Inefficient management of overdue books:} Tracking due dates and identifying overdue books requires manual checking, which may result in delays and inaccuracies.

\item \textbf{Risk of data loss:} Physical records can be damaged, misplaced, or lost due to accidents or improper handling.

\item \textbf{Limited accessibility:} Members must visit the library physically to check book availability or obtain information about library resources.

\item \textbf{Scalability issues:} Managing a large collection of books and a growing number of members becomes increasingly difficult using manual methods.

\end{itemize}

\subsection{Proposed System}
The proposed LMS is a web application with a relational database. It automates catalogue management, enforces borrowing rules, provides instant search, and allows online reservations. It reduces manual errors and improves accessibility.


\section{Software Process Model}

\subsection{Waterfall Model}
The Waterfall model (classic life cycle) is chosen: Requirements $\rightarrow$ Design $\rightarrow$ Implementation $\rightarrow$ Testing $\rightarrow$ Deployment $\rightarrow$ Maintenance.

\subsection{Process Model Diagram}
\textbf{Include necessary diagrams}

\subsection{Phases of the Model}
\begin{enumerate}[leftmargin=*]
  \item \textbf{Requirements analysis:} This SRS is the output.
  \item \textbf{System design:} Architecture, database schema, component design.
  \item \textbf{Implementation:} Coding the web application.
  \item \textbf{Testing:} Unit, integration, and system testing against the SRS.
  \item \textbf{Deployment:} Install on library server.
  \item \textbf{Maintenance:} Bug fixes and minor enhancements.
\end{enumerate}

\subsection{Justification}
\begin{itemize}[leftmargin=*]
  \item Requirements of LMS are well‑understood and stable (typical library operations).
  \item The project is small to medium in size, with a clear sequence of phases.
  \item The team has previous experience with similar academic projects.
  \item Milestones (SRS, design document, tested system) are easy to evaluate.
\end{itemize}


\section{Overall Description}

\subsection{Product Perspective}
% Is it standalone or part of a larger system? Describe its context.
The \projecttitle{} is a standalone web application backed by a relational
database. It interacts with library staff and members through a browser.

\subsection{Product Functions}
% A short bulleted summary of the major functions (detailed in Section 3).
\begin{itemize}[leftmargin=*]
  \item Manage the book catalogue (add / edit / remove books).
  \item Search the catalogue by title or author.
  \item Issue and return books, enforcing borrowing limits and due dates.
  \item Reserve unavailable books.
  \item Manage member records.
\end{itemize}

\subsection{User Characteristics}
\begin{itemize}[leftmargin=*]
  \item \textbf{Librarian:} trained staff member; performs all catalogue and loan operations.
  \item \textbf{Member:} general user; searches and reserves books.
  \item \textbf{System Administrator:} Maintains servers, backups, and user accounts.
\end{itemize}

\subsection{Key Entities}
\begin{itemize}[leftmargin=*]
  \item \textbf{Book:} ISBN, title, author, publisher, year, quantity, available copies.
  \item \textbf{Member:} Member ID, name, contact, maximum allowed books.
  \item \textbf{Transaction:} Issue date, due date, return date, book copy, member.
  \item \textbf{Reservation:} Book ID, member ID, reservation date, status.
\end{itemize}

\subsection{Constraints}
% Technology, regulatory, hardware, or institutional limits.
\begin{itemize}[leftmargin=*]
  \item Must run on the institution's existing server infrastructure.
  \item Passwords must be stored only as salted hashes.
\end{itemize}

\subsection{Assumptions and Dependencies}
\begin{itemize}[leftmargin=*]
  \item Members are registered before they can reserve books.
  \item A network connection to the database server is available.
\end{itemize}

%==============================================================
%  3. SPECIFIC REQUIREMENTS
%==============================================================
\section{Specific Requirements}

\subsection{External Interface Requirements}
\subsubsection{User Interfaces}
A browser-based interface with a catalogue search page, a librarian dashboard,
and a member account page.
\subsubsection{Hardware Interfaces}
Standard PC / server hardware; optionally a barcode scanner for ISBNs.
\subsubsection{Software Interfaces}
A relational database (e.g., PostgreSQL) accessed by the application.
\subsubsection{Communication Interfaces}
HTTPS for all client-server traffic.

\subsection{Functional Requirements}
% Each requirement: an ID, a single testable "shall" statement, and a priority.
% Copy a row to add more. Keep IDs unique and sequential.
\renewcommand{\arraystretch}{1.35}
\begin{tabularx}{\textwidth}{|c|X|c|}
\hline
\reqheader{ID}{Functional Requirement}{Priority}
FR-1 & The system shall allow a librarian to add, edit, and remove books in the catalogue. & High \\ \hline
FR-2 & The system shall allow a member to search the catalogue by title or author. & High \\ \hline
FR-3 & The system shall allow a librarian to issue an available book to a registered member. & High \\ \hline
FR-4 & The system shall reject an issue when the member already holds the maximum allowed books. & High \\ \hline
FR-5 & The system shall set a due date 14 days from the issue date. & Medium \\ \hline
FR-6 & The system shall allow a member to reserve a book that is currently unavailable. & Low \\ \hline
% FR-7 & [ your requirement ] & [ High/Med/Low ] \\ \hline
\end{tabularx}

\subsection{Non-functional Requirements}
\renewcommand{\arraystretch}{1.35}
\begin{tabularx}{\textwidth}{|c|X|c|}
\hline
\reqheader{ID}{Non-functional Requirement}{Category}
NFR-1 & A catalogue search shall return results within 2 seconds for a 10,000-book catalogue. & Performance \\ \hline
NFR-2 & Passwords shall be stored only as salted hashes. & Security \\ \hline
NFR-3 & The system shall be usable on a 1366$\times$768 screen without horizontal scrolling. & Usability \\ \hline
NFR-4 & The system shall be available 99\% of the time during library opening hours. & Reliability \\ \hline
% NFR-5 & [ your requirement ] & [ category ] \\ \hline
\end{tabularx}

\subsection{Other Requirements}
% Anything that doesn't fit above (e.g., data retention, legal, localisation).
If any.

%==============================================================
%  4. CONCLUSION  (lab-report wrap-up)
%==============================================================
\section{Conclusion}
% 3--5 sentences: what you produced, what you learned, any difficulties.
In this lab the team elicited and documented the requirements of the
\projecttitle{} using the IEEE~830 template. The functional and non-functional
requirements were each given unique identifiers so they can be traced into the
design and test phases in later labs.

\section{References}
\begin{enumerate}[leftmargin=*]
  \item IEEE Std 830-1998, Recommended Practice for Software Requirements Specifications.
  \item more reference if any.
\end{enumerate}
\end{document}